% Appendix

\section{Feature Extraction Details}
\label{app:features}

The 64-dimensional feature vector used by the MLP classifier is constructed as follows:

\begin{table}[h]
\centering
\caption{Feature vector composition.}
\label{tab:features}
\begin{tabular}{clc}
\toprule
\textbf{Index} & \textbf{Feature} & \textbf{Dims} \\
\midrule
0 & Normalized syscall number & 1 \\
1--10 & Syscall category (one-hot) & 10 \\
11--20 & Dangerous syscall indicators & 10 \\
21 & Is root user (uid=0) & 1 \\
22 & Privilege escalation flag & 1 \\
23--32 & N-gram frequency statistics & 10 \\
33--42 & Sliding window syscall counts & 10 \\
43--52 & Parent process features & 10 \\
53--62 & Temporal features & 10 \\
63 & STIDE anomaly score & 1 \\
\bottomrule
\end{tabular}
\end{table}

\textbf{Syscall categories} group related syscalls:
\begin{itemize}
    \item Process: execve, fork, clone, exit
    \item File: open, read, write, close, unlink
    \item Network: socket, connect, accept, sendto, recvfrom
    \item Memory: mmap, mprotect, brk
    \item IPC: pipe, shmget, msgget
    \item System: ioctl, ptrace, mount
\end{itemize}

\textbf{Dangerous syscall indicators} flag potentially suspicious syscalls:
\begin{itemize}
    \item ptrace (debugging/injection)
    \item execve from unusual paths (/tmp, /dev/shm)
    \item setuid/setgid (privilege changes)
    \item module\_init (kernel module loading)
    \item mount with unusual options
\end{itemize}

\section{Training Hyperparameters}
\label{app:hyperparams}

\begin{table}[h]
\centering
\caption{Default training hyperparameters.}
\label{tab:hyperparams-app}
\begin{tabular}{lc}
\toprule
\textbf{Parameter} & \textbf{Value} \\
\midrule
Learning rate & 0.01 \\
Epochs & 200 \\
Batch size & 64 \\
Hidden layers & [128, 64, 32, 16] \\
Activation (hidden) & ReLU \\
Activation (output) & Sigmoid \\
Loss function & Weighted BCE \\
Optimizer & SGD with momentum (0.9) \\
Weight decay & $10^{-4}$ \\
Gradient clipping & 1.0 \\
Early stopping patience & 20 epochs \\
\bottomrule
\end{tabular}
\end{table}

\section{ADFA-LD Attack Categories}
\label{app:adfa}

The ADFA-LD dataset contains the following attack types:

\begin{table}[h]
\centering
\caption{ADFA-LD attack categories.}
\label{tab:adfa-attacks}
\begin{tabular}{lcp{4cm}}
\toprule
\textbf{Category} & \textbf{Traces} & \textbf{Description} \\
\midrule
Hydra FTP & 162 & Brute force FTP login \\
Hydra SSH & 148 & Brute force SSH login \\
Adduser & 91 & Unauthorized user creation \\
Java Meterpreter & 125 & Remote shell via Meterpreter \\
Linux Meterpreter & 98 & Native Meterpreter payload \\
Web shell & 122 & PHP/Python web backdoors \\
\midrule
Total attacks & 746 & \\
Normal & 5,205 & Standard system operation \\
\bottomrule
\end{tabular}
\end{table}

\section{Detailed Latency Breakdown}
\label{app:latency}

\begin{table}[h]
\centering
\caption{Per-component latency breakdown.}
\label{tab:latency-breakdown}
\begin{tabular}{lc}
\toprule
\textbf{Component} & \textbf{Latency} \\
\midrule
Event read from ring buffer & 0.2$\mu$s \\
Feature extraction & 3.1$\mu$s \\
STIDE n-gram check & 0.2$\mu$s \\
MLP forward pass & 12.8$\mu$s \\
Result formatting & 0.2$\mu$s \\
\midrule
\textbf{Total (with MLP)} & \textbf{16.5$\mu$s} \\
\textbf{Total (STIDE only)} & \textbf{3.5$\mu$s} \\
\bottomrule
\end{tabular}
\end{table}

The MLP forward pass dominates latency (78\%).
Events filtered by STIDE avoid this cost, achieving 3.5$\mu$s latency.
